\chapter{Fallopian Tube Infection}

\section*{Definition and Overview}
A fallopian tube infection is an infection of one or both of the fallopian tubes, the tubes that carry eggs from the ovaries to the uterus. Fallopian tube infection is part of pelvic inflammatory disease (PID), which also affects the uterus, cervix, and ovaries. It is most often caused by sexually transmitted infections, particularly chlamydia and gonorrhoea, that spread upwards from the cervix. If untreated, the infection can scar the tubes, causing chronic pain, ectopic pregnancy, and infertility. Prompt treatment with antibiotics prevents most complications.

\section*{Epidemiology}
Pelvic inflammatory disease affects a significant number of women of reproductive age each year, and it is most common in sexually active women under 25. Chlamydia and gonorrhoea are the main causes, and many women with these infections have no symptoms, allowing silent spread. The incidence has declined with screening and treatment of STIs, but it remains an important cause of preventable infertility.

\section*{Aetiology and Pathophysiology}
Fallopian tube infection is caused by bacteria ascending from the lower genital tract.
\begin{itemize}
    \item \textbf{STIs:} chlamydia and gonorrhoea are the most common causes
    \item \textbf{Other bacteria:} organisms from the vagina and bowel, including mycoplasma and anaerobes
    \item \textbf{Spread:} bacteria travel from the cervix through the uterus to the tubes and ovaries
    \item \textbf{Risk factors:} young age, multiple sexual partners, a new partner, previous STIs, and an intrauterine device
    \item \textbf{Mechanism:} the bacteria cause inflammation of the tubes, which swell and may fill with pus
    \item \textbf{Scarring:} healing leaves scar tissue that narrows or blocks the tubes
    \item \textbf{Consequences:} blockage prevents the egg and sperm from meeting, and scarring increases the risk of ectopic pregnancy
\end{itemize}

\section*{Clinical Features}
Symptoms range from none to severe.
\begin{itemize}
    \item Lower abdominal pain, which may be on one or both sides
    \item Abnormal vaginal discharge
    \item Pain during intercourse
    \item Bleeding between periods or after sex
    \item Fever and feeling unwell in severe cases
    \item Urinary symptoms, including pain on urination
    \item Many women have no symptoms, and the damage is found only when fertility is affected
\end{itemize}

\section*{Differential Diagnosis}
Other conditions cause similar symptoms.
\begin{itemize}
    \item Ectopic pregnancy, which is an emergency
    \item Ovarian cysts, including torsion
    \item Appendicitis
    \item Endometriosis
    \item Urinary tract infection
    \item Gastroenteritis and other causes of abdominal pain
\end{itemize}

\section*{When to Seek Medical Care}
Any woman with lower abdominal pain and discharge, or with a risk of STIs and pelvic symptoms, should be assessed by a doctor promptly, because early treatment prevents scarring.

\textbf{Call 000 (or local emergency services) for:}
\begin{itemize}
    \item Sudden severe abdominal or pelvic pain
    \item Fever with pelvic pain
    \item Pain with vomiting or collapse
    \item Shoulder tip pain with pelvic pain, which may indicate an ectopic pregnancy
    \item Foul-smelling discharge with severe illness
\end{itemize}

\section*{Diagnosis and Investigations}
Diagnosis is clinical, with tests to confirm.
\begin{itemize}
    \item \textbf{History:} symptoms, sexual history, and risk factors
    \item \textbf{Pelvic examination:} tenderness of the uterus and tubes, and discharge
    \item \textbf{Swabs:} tests for chlamydia, gonorrhoea, and other organisms
    \item \textbf{Urine tests:} for STIs and urinary infection
    \item \textbf{Blood tests:} inflammatory markers and pregnancy test
    \item \textbf{Ultrasound:} to assess the tubes, ovaries, and any abscess
\end{itemize}

\section*{Management}
Treatment is with antibiotics, and prompt treatment protects fertility.
\begin{itemize}
    \item \textbf{Antibiotics:} a course that covers the likely organisms; severe infection is treated in hospital with intravenous antibiotics
    \item \textbf{Partner treatment:} sexual partners are tested and treated to prevent reinfection
    \item \textbf{Pain relief:} anti-inflammatory medicines
    \item \textbf{Rest:} avoiding sex and strenuous activity until symptoms settle
    \item \textbf{Follow-up:} confirmation that the infection has cleared
    \item \textbf{Abscess:} drainage or surgery for a tubo-ovarian abscess
    \item \textbf{Screening:} testing for other STIs
\end{itemize}

\section*{Complications}
\begin{itemize}
    \item Scarring and blockage of the tubes, causing infertility
    \item An increased risk of ectopic pregnancy, which can be life-threatening
    \item Chronic pelvic pain
    \item Tubo-ovarian abscess
    \item Spread of infection, including peritonitis and sepsis
    \item Recurrent infection
\end{itemize}

\section*{Prognosis and Outcomes}
With early antibiotic treatment, most women recover fully with no lasting damage. The risk of infertility rises with each episode of infection, which is why prompt treatment, partner treatment, and STI screening are so important. Women who have had PID should be aware of their higher risk of ectopic pregnancy in future pregnancies.

\section*{Prevention and Self-Care}
\begin{itemize}
    \item Practise safe sex and use condoms
    \item Get tested for STIs regularly, especially with a new partner
    \item Complete the full course of antibiotics
    \item Ensure sexual partners are treated
    \item Avoid sex until the infection has cleared
    \item Attend follow-up to confirm the infection is gone
    \item Seek early care for pelvic symptoms
\end{itemize}

\section*{Sources and Further Reading}
The following reputable sources provide further information for healthcare students and patients seeking reliable, current advice about fallopian tube infection and PID.
\begin{itemize}
    \item Jean Hailes for Women's Health. \textit{Pelvic inflammatory disease.} https://www.jeanhailes.org.au/
    \item Healthdirect Australia. \textit{Pelvic inflammatory disease.} https://www.healthdirect.gov.au/pelvic-inflammatory-disease
    \item Family Planning Australia. \textit{STI testing and information.} https://www.fpnsw.org.au/
    \item NHS. \textit{Pelvic inflammatory disease.} https://www.nhs.uk/conditions/pelvic-inflammatory-disease/
    \item Mayo Clinic. \textit{Pelvic inflammatory disease.} https://www.mayoclinic.org/diseases-conditions/pelvic-inflammatory-disease/
    \item MSD Manual. \textit{Pelvic inflammatory disease.} https://www.msdmanuals.com/
\end{itemize}
